\chapter{Research Method}\label{ch:research-method}

\section{Research Design}\label{sec:research-design}

The study combined artifact development with a controlled empirical evaluation. The developed artifact was the Guided Intelligence Workspace retrieval pipeline, whose purpose was to construct source-grounded evidence for explaining repository-level issues. Formative development established the final pipeline design, after which the evaluation assessed its retrieval quality, efficiency, and stability under the frozen comparison described below.

Artifact development and final evaluation were treated as separate phases. During formative development, the development partition was used to expose retrieval failures and assess proposed changes. A change began with an observed limitation and a bounded hypothesis about the responsible pipeline behaviour. It was first checked with focused tests and diagnostic executions, followed by repeated executions of the actual retrieval pipeline when the change affected retrieval behaviour. Changes were retained when their results were sufficiently repeatable and did not introduce unacceptable regressions; otherwise, they were revised or removed. These experiments informed the final architecture but were not included in the reported five-condition evaluation. The principal retained and rejected design experiments are presented in the Design Rationale chapter because their purpose was to explain architectural decisions rather than estimate final system performance.

After formative development, the system variants and evaluation procedure were frozen for the final campaign. Results were reported separately for the development and held-out partitions before a combined summary, distinguishing behaviour on cases that had informed development from behaviour on unseen cases.

\section{CodeRepoQA Corpus and Historical Case Construction}\label{sec:coderepoqa-corpus-and-historical-case-construction}

The testcase corpus was motivated by CodeRepoQA, introduced by Hu et al. in \emph{Understanding Large Language Model Performance in Software Engineering: A Large-scale Question Answering Benchmark} \cite{Hu2025CodeRepoQA}. CodeRepoQA uses multi-turn conversations from GitHub issues to evaluate repository-level question answering. This study used issues from its released TypeScript, pandas, and Vue collections as the initial candidate pool, but did not adopt its original dialogue-answering protocol. The issue categories, case-selection rules, pre-resolution snapshots, evaluator Oracles, and development and held-out partitions described below were constructed specifically for this retrieval evaluation.

The resulting corpus contained 35 retrieval-grounded cases: 11 from TypeScript, 12 from pandas, and 12 from Vue. These cases were distributed evenly across seven issue categories: bug and regression, feature enhancement, performance and memory, compatibility and versioning, API behaviour and design, testing and build tooling, and maintenance and refactoring. Each category contained five cases. Three additional usage-question cases were retained for explanation-oriented analysis but excluded from file-ranking comparisons because they did not provide equivalent implementation-file Oracles.

Cases were selected only when the initial issue described a concrete repository-level problem and a fixing pull request, commit, or comparable resolution artifact could be identified. Selection also favoured issues whose problem could be understood without later discussion, whose resolution provided sufficient evidence to construct an Oracle, and whose responsibility did not duplicate another selected case. The resulting corpus included both compact single-file changes and mechanisms spanning implementation, tests, documentation, or multiple subsystems. This variation allowed retrieval to be compared across different issue intents and repository structures rather than on repeated instances of one narrow task.

The 35 retrieval-grounded cases were divided into 28 development cases and seven final cases. The development partition contained four cases from every issue category and was available during formative retrieval experiments. The final partition contained one held-out case from each category and was not used to tune retrieval behaviour, providing a check on whether observed behaviour extended beyond the cases that informed system development.

Each case paired the visible issue with a repository snapshot from before its resolution. Where verified pull-request metadata provided a base commit, that commit defined the snapshot. Otherwise, the harness used the parent of an identified fixing or event commit, or, with lower confidence, the latest available commit preceding the issue's creation time. The selected commit was materialised as a detached checkout and reused across conditions. Retrieval received only the issue title, creation time, body, and the pre-resolution repository. Later comments, issue events, fixing references, pull-request metadata, diffs, and \texttt{verification.json} remained evaluator-only material. This separation prevented the systems from retrieving the known fix or using later maintainer explanations as if they had been available when the issue was opened.

Index scope was defined independently of each case's Oracle. Common exclusions removed version-control data, dependency directories, virtual environments, caches, coverage data, build output, and retrieval-system state. For TypeScript, generated compiler-test baselines and the entire \texttt{lib} directory were additionally excluded because the large generated compiler and server bundles made unscoped indexing impractical, while authored compiler test inputs under \texttt{tests/cases} remained searchable. Excluding \texttt{lib} also removed declaration files and therefore limited the declaration-oriented evidence available to every condition. The same effective exclusion list was supplied to BM25, Qdrant, and CodeGraph indexing so that each stage received the same eligible repository paths. Repository snapshots and exclusions were then held constant across the evaluated conditions; no path was removed because it was known to be irrelevant to an individual Oracle.

\section{Oracle Construction}\label{sec:oracle-construction}

Each retrieval-grounded testcase had a hidden evaluator Oracle that was stored separately from the issue presented to the evaluated system. The Oracle was constructed from post-resolution artifacts, principally the changed files of a locally available fixing commit or merged pull request. Issue events, linked changes, and maintainer discussion provided additional resolution context. These sources were used only to construct and verify the evaluator metadata. The resulting Oracle was frozen in the testcase's \texttt{verification.json} before the measured runs so that no condition could influence its definition.

Oracle files were divided according to their role in the resolution. Implementation files contained the production behavior changed to resolve the issue. Test or validation files exercised that behavior or encoded the regression, while documentation files recorded a relevant user-facing or maintenance consequence. The latter two categories were treated as supporting evidence rather than substitutes for the responsible implementation. Every case included in file-ranking statistics was required to contain at least one implementation file.

Oracle and retrieved files were matched by exact repository-relative path after normalization and duplicate removal. For file-level ranking measures, the ordered final evidence list was collapsed into unique files. When multiple evidence items came from the same file, that file inherited the position of its highest-ranked item. The implementation and supporting categories supplied the relevance assignments defined in the evaluation-measures subsection.

Final evidence selection also produced \texttt{coverage\_status} and \texttt{sufficient} outputs. These values represented the selector's own judgement of whether the retrieved evidence was complete or partial; because the corpus did not provide an independent, uniform line-level or structural-owner-level completeness Oracle, they were treated as system outputs rather than ground-truth labels. Fixing changes likewise provided an imperfect file Oracle: a pull request could contain incidental files, omit relevant unchanged context, or address more than one concern. Separating implementation from supporting files and retaining the hidden resolution context reduced these risks but did not eliminate them.

\section{Evaluated Systems and Comparison Design}\label{sec:evaluated-systems-and-comparison-design}

The empirical evaluation compared five repository-retrieval conditions: the complete Guided Intelligence Workspace pipeline, Workspace without CodeGraph, Workspace without its adaptive controller, Workspace without either CodeGraph or the adaptive controller, and Codex retrieval. The four Workspace conditions formed a two-factor design in which structural graph support and adaptive exploration were independently enabled or disabled. Codex retrieval formed a broader complete-system comparison rather than an ablation of Workspace.

\begingroup
\small
\begin{longtable}{@{}L{0.28\linewidth}L{0.12\linewidth}L{0.16\linewidth}L{0.34\linewidth}@{}}
\caption{Evaluated retrieval conditions}\label{tab:evaluated-retrieval-conditions-1}\\
\toprule
\textbf{Condition} & \textbf{CodeGraph} & \textbf{Adaptive controller} & \textbf{Experimental role} \\
\midrule
\endfirsthead
\multicolumn{4}{c}{\tablename\ \thetable\ -- continued}\\
\toprule
\textbf{Condition} & \textbf{CodeGraph} & \textbf{Adaptive controller} & \textbf{Experimental role} \\
\midrule
\endhead
\midrule
\multicolumn{4}{r}{Continued on next page}\\
\endfoot
\bottomrule
\endlastfoot
Full Workspace & Enabled & Enabled & Reference native pipeline \\
Workspace without CodeGraph & Disabled & Enabled & Structural-retrieval ablation \\
Workspace without adaptive controller & Enabled & Disabled & Adaptive-exploration ablation \\
Workspace without CodeGraph and without adaptive controller & Disabled & Disabled & Combined ablation and interaction condition \\
Codex retrieval & Not applicable & Not applicable & External complete-system comparison \\
\end{longtable}
\endgroup

Full Workspace served as the native reference condition and enabled both CodeGraph-based structural processing and post-round-zero adaptive exploration. It followed three principal phases: initial evidence construction, bounded adaptive completion, and final evidence consolidation. The internal architecture of these phases is detailed in Chapter 5. During the evaluation runs, response generation was disabled, while final evidence selection remained enabled because each condition could change the evidence available to the selection model and thereby change its decision. Retaining this stage also preserved the complete retrieval flow used by the system rather than evaluating only an intermediate candidate pool.

The first native ablation set \texttt{structural\_graph\_enabled=false}. This disabled CodeGraph-based owner resolution, graph evidence, and graph-dependent navigation; structural operations returned explicit empty results. Qdrant's lexical and semantic retrieval channels remained enabled, and retrieved ranges remained available as range-level observations for semantic qualification and adaptive exploration. Final evidence selection also remained unchanged. This contrast isolated the contribution of CodeGraph-based structural grounding.

The second native ablation set \texttt{adaptive\_controller\_enabled=false}. This disabled adaptive exploration after round zero, including controller-directed follow-up retrieval and qualification. Initial retrieval and round-zero semantic qualification remained enabled, and the resulting evidence passed directly to the unchanged final-consolidation stages. This contrast isolated the contribution of post-round-zero controller exploration.

The combined native ablation set both configuration flags to \texttt{false}. It neither resolved retrieved ranges through CodeGraph nor performed adaptive exploration after round zero. The remaining lexical and semantic evidence proceeded through the unchanged final-consolidation boundary. This condition completed the factorial design and made it possible to determine whether the contribution of either capability depended on the presence of the other.

Codex retrieval provided the external agentic comparison. It used \texttt{gpt-5.6-luna} with the frozen \texttt{efficient} prompt profile and inspected the pre-resolution repository through its own retrieval process rather than the Workspace stages. Its evidence was evaluated against the same frozen Oracles. Differences could arise from search, source inspection, context construction, navigation, stopping, or consolidation; the comparison therefore characterised complete-system behaviour and did not isolate a single component.

Across the four Workspace conditions, the initial comparison and qualification contracts, final selector, and run procedure remained fixed. Only the availability of CodeGraph and post-round-zero adaptive exploration differed according to the two feature flags; the shared input, snapshot, and index controls were those defined in the corpus-construction subsection.

The factorial design supported four controlled component contrasts. Full Workspace versus Workspace without CodeGraph measured the effect of removing CodeGraph while adaptive exploration remained available. Workspace without the adaptive controller versus the combined ablation measured the same structural effect without adaptive exploration. Full Workspace versus Workspace without the adaptive controller measured the effect of removing adaptive exploration while CodeGraph remained available. Workspace without CodeGraph versus the combined ablation measured the controller effect without structural graph support. Comparing these paired effects established whether the contribution of CodeGraph changed with controller availability, and whether adaptive exploration depended on CodeGraph-derived owners and relationships. Full Workspace versus Codex remained a separate complete-system comparison between the native evidence lifecycle and an external agentic retrieval process.

\section{Quantitative Evaluation Measures}\label{sec:quantitative-evaluation-measures}

Relevant-file retrieval and ranking were measured at cutoffs of 1, 2, 5, and 10 ordered unique files. Precision at rank \emph{k} (P@\emph{k}) was the number of implementation Oracle files in the first \emph{k} positions divided by \emph{k}. The denominator remained \emph{k} when fewer files were returned, with unfilled positions treated as non-relevant. Recall at rank \emph{k} (R@\emph{k}) was the number of implementation Oracle files in the first \emph{k} positions divided by the total number of implementation Oracle files for that case. Precision therefore measured the concentration of responsible files near the top of the result, whereas recall measured recovery of the known implementation set.

Normalised discounted cumulative gain (NDCG) was calculated at the same cutoffs to account for rank and graded relevance. Implementation files received relevance grade 2, supporting test, validation, or documentation files grade 1, and all other files grade 0. At rank \emph{i}, discounted gain was calculated as $(2^{grade}-1)/\log_2(i+1)$. The sum through rank \emph{k} was divided by the ideal discounted gain for the same case, with grade-2 files placed before grade-1 files. NDCG consequently ranged from zero to one and rewarded early implementation evidence most strongly while still recognising supporting evidence.

The quantitative comparison also recorded the number of selected evidence items and ordered unique files returned by each run. These counts described the size of the final result and helped distinguish ranking differences from differences in how many files a condition returned. The cross-condition efficiency comparison used end-to-end runtime and provider-reported retrieval-flow token usage. Response-generation tokens were outside the evaluation, and indexing usage was not included in the reported flow-token totals.

Because the LLM-backed stages remained potentially stochastic, each case-condition pair had four valid repetitions under a fixed configuration. Across five conditions and 35 retrieval-grounded cases, this produced 700 valid runs. Run-to-run variation was described using the frequency of any implementation-file hit and full implementation recall, the mean pairwise Jaccard similarity of retrieved-file sets, and the population standard deviation of R@10. For every ranking and operational measure, values were first averaged across the four repetitions of one case and then macro-averaged across case means. This gave every case equal weight.

A run was valid only when the intended condition and repository snapshot completed with parseable artifacts, a non-failed coverage state, and non-empty usable evidence. Infrastructure interruptions, unavailable services, schema-invalid outputs, incomplete artifacts, and zero-evidence executions remained recorded in the campaign ledgers but were excluded from aggregates. Collection continued until four valid repetitions were available for every case and condition. Valid runs were not replaced based on their scores, maintaining an auditable separation between execution failure handling and retrieval quality.

\section{Reproducibility, Validity, and Ethical Scope}\label{sec:reproducibility-validity-and-ethical-scope}

Each execution produced a timestamped run ID and a directory containing the run metadata, retrieval result, evaluator comparison, and available stage records. The metadata identified the testcase, pre-resolution repository commit, effective source exclusions, retrieval mode, relevant feature flags, model configuration, and whether response generation and final evidence selection were enabled. Versioned configuration files defined the five evaluated conditions, while campaign ledgers recorded every accepted run and failed attempt together with its timestamps and failure reason. The final aggregation linked each reported observation to its run ID and was generated by a repository-local analysis script. These records preserved the connection from aggregate values to individual execution artifacts without placing the complete run inventory in the main thesis narrative.

Index reuse was controlled by the repository snapshot and indexing configuration. BM25 and Qdrant indexes were constructed over the same eligible paths and reused across the four Workspace conditions; the ablations changed the availability of CodeGraph and adaptive exploration rather than the underlying lexical or vector corpus. Repeated runs also reused the prepared snapshot and indexes instead of rebuilding them. This removed index construction as a source of within-case variation, although it means that reported runtime and token totals describe retrieval execution rather than the cost of preparing a repository for first use.

The retained artifacts permit the reported runs and calculations to be audited, but they do not guarantee bit-for-bit reproduction: hosted model behaviour, service infrastructure, and the Codex execution environment can change independently of the repository configuration. Runtime additionally depends on machine load, network latency, and provider response time and was interpreted as an operational measurement of this campaign rather than a platform-independent constant.

Several validity boundaries constrained the interpretation. The controlled factorial design supported attribution to the two Workspace capabilities, whereas the Workspace-Codex contrast supported only complete-system conclusions. Construct validity was bounded by the file-based Oracles and self-assessed completeness outputs described above: the ranking measures established recovery and ordering of known fixing files, not whether the evidence formed a complete causal explanation. External validity was limited by the historical corpus, the three selected open-source repositories, their language ecosystems and project structures, and the repository-specific exclusions.

The empirical material consisted of publicly available open-source repositories, issue reports, and resolution artifacts; the study involved no recruited participants or intervention in repository communities. Author identities were not used as study variables, and the evaluation did not require private repository content or confidential research data to be submitted to model providers. Any generative AI output used during the research was treated as unverified assistance, while responsibility for testcase selection, Oracle construction, experimental decisions, source verification, and analysis remained with the researcher. The broader risks of unsupported repository explanations, privacy and intellectual-property exposure, and overreliance on apparently plausible output are examined in the Discussion.
